\SourceLeaf{260}
\clearpage
\Article{9}
All this had already been treated in almost the same way in the Disquis.\ Ar.; some things, however, had been added there with their proof omitted, and we think it appropriate to supply that proof here. We stated there that finding the value of the radical quantity ${}^{\beta}\!\surd\theta'$, which, since $\theta'$ is an imaginary quantity, seems to require dividing both a ratio and an angle into $\beta$ parts, depends on the latter alone, and that the former can always be reduced to extracting a single square root. We prove this as follows.

Denoting the imaginary quantity $\surd-1$ by $i$, as above, and writing $\theta'=P+iQ$ and ${}^{\beta}\!\surd\theta'=p+iq$, using any one value of the radical and taking $P,Q,p,q$ to be real, it is known that, if the positive quantities $E,e$ and the angles $F,f$ are determined so that $P=E\cos F$, $Q=E\sin F$, $p=e\cos f$, $q=e\sin f$, then $e={}^{\beta}\!\surd E$, and $f$ equals one of the angles
\[
\frac1\beta F,\quad\frac1\beta(F+360^0),\quad\frac1\beta(F+720^0)\dotsfour\frac1\beta(F+(\beta-1)360^0)
\]
Thus $f$\textsuperscript{.}\ ger is determined by division of the angle $F$ into $\beta$ parts, but we can avoid extracting the root ${}^{\beta}\!\surd E$ as follows. Every product $r^kR^K$ has the same real part as $r^{-k}R^{-K}$, while the imaginary parts containing the factor $i$ in these products are equal and opposite. It follows at once that $[r^{-1},R^{-1}]=p-iq=e(\cos f-i\sin f)$, and hence
\[
[r,R]\times[r^{-1},R^{-1}]=e^2
\]

But by the theorem of Article~4 this product becomes
\[
\begin{aligned}
&{}=[1,1]+R[r^{g^\alpha-1},1]+R^2[r^{g^{2\alpha}-1},1]+\text{etc.}\\
&\qquad{}+R^{\beta\gamma-1}[r^{g^{\alpha\beta\gamma-\alpha}-1},1]\\
&{}=\beta\gamma+R(\beta\gamma,g^\alpha-1)+R^2(\beta\gamma,g^{2\alpha}-1)+\text{etc.}\\
&\qquad{}+R^{\beta\gamma-1}(\beta\gamma,g^{\alpha\beta\gamma-\alpha}-1)
\end{aligned}
\]
This quantity can be determined if $R$ and all the periods of $\beta\gamma$ terms are assumed known. Determining $e$ therefore requires only the extraction of a square root.

In the special case $\alpha=1$, the individual periods $(\beta\gamma,g^\alpha-1)$, $(\beta\gamma,g^{2\alpha}-1)$, etc., are evidently $=r+r^2+r^3+r^4+\text{etc.}+r^{n-1}$, and hence
\[
\begin{aligned}
ee&=\beta\gamma+(R+R^2+R^3+\text{etc.}+R^{\beta\gamma-1})(r+r^2+r^3+\text{etc.}+r^{n-1})\\
&=\beta\gamma+1=n
\end{aligned}
\]
\SourceLeaf{261}
provided that $r$ and $R$ are assumed to denote roots different from $1$, and therefore we always have $e=\surd n$ (Disq. arithm., end of Article~360).

\Needspace{19\baselineskip}
\Article{10}
So far we have conducted our investigation in the greatest generality, to cover all values of the numbers $\alpha,\beta,\gamma$; from now on, however, we shall pass to the more restricted case $\alpha=1$, which will open the way to most fruitful and elegant investigations. Thus the symbol $[r,R]$ will denote the function
\[
r+Rr^g+R^2r^{g^2}+R^3r^{g^3}+\text{etc.}+R^{n-2}r^{g^{n-2}}
\]
where $n$ is a prime number, $r$ an arbitrary root of the equation $x^n-1=0$ (the root $1$ not excluded), $R$ an arbitrary root of the equation $x^\beta-1=0$, with $\beta$ denoting a given divisor of $n-1$, and finally $g$ an integer that is a fixed primitive root modulo $n$. For brevity we shall also write
\[
\begin{aligned}
1+r+r^2+r^3+\text{etc.}+r^{n-1}&=s\\
1+R+R^2+R^3+\text{etc.}+R^{n-2}&=S
\end{aligned}
\]
from which it is evident that $s$ becomes $=n$ for $r=1$, but $s=0$ for every other value of $r$, and similarly $S=n-1$ for $R=1$, but $S=0$ for every other value of $R$.

By Article~3 we therefore have $[1,R]=S$, $[r,1]=s-1$; further, for every integer value of $m$ not divisible by $n$, $[r^m,R]=R^{-\operatorname{ind}m}[r,R]$, or more generally $[r^m,R^M]=R^{-M\operatorname{ind}m}[r,R^M]$, where $\operatorname{ind}m$ is the exponent of the power of $g$ congruent modulo $n$ to $m$. Applying this transformation to what we showed in Article~5, we find that a product of two or more functions such as $[r^h,R^H]$ reduces to the form
\[
A[1,R^\lambda]+B[r,R^\lambda]
\]
where $A$ and $B$ are rational functions of $R$ with integer coefficients, and $\lambda$ is the sum of all the values of $H$. It will be of great importance to put such transformations into an efficient algorithm; for this purpose we must first consider more closely the nature of the product of two functions.

\clearpage
\Article{11}
The product $[r,R^\mu]\times[r,R^\nu]$ becomes, by the theorem of Article~4, $=$
\SourceLeaf{262}
\[
\begin{aligned}
&[r^2,R^{\mu+\nu}]+R^\mu[r^{g+1},R^{\mu+\nu}]+R^{2\mu}[r^{g^2+1},R^{\mu+\nu}]\\
&\qquad+R^{3\mu}[r^{g^3+1},R^{\mu+\nu}]+\text{etc.}+R^{(n-2)\mu}[r^{g^{n-2}+1},R^{\mu+\nu}]
\end{aligned}
\]
Among the $n-1$ exponents $2,g+1,g^2+1,g^3+1$, etc., $g^{n-2}+1$, exactly one is divisible by $n$, namely $g^{\frac12(n-1)}+1$, so that the corresponding term of our sum will be $R^{\frac12(n-1)\mu}[1,R^{\mu+\nu}]$. This term will be $=0$ whenever we do not have $R^{\mu+\nu}=1$, and $=(n-1)R^{\frac12(n-1)\mu}=\pm(n-1)$ for $R^{\mu+\nu}=1$. The remaining parts of our sum, whose sum we shall put $=\Omega$, are transformed as follows:
\[
\begin{aligned}
\relax[r^2,R^{\mu+\nu}]&=R^{-(\mu+\nu)\operatorname{ind}2}[r,R^{\mu+\nu}]\\
R^\mu[r^{g+1},R^{\mu+\nu}]&=R^{\mu-(\mu+\nu)\operatorname{ind}(g+1)}[r,R^{\mu+\nu}]\\
R^{2\mu}[r^{g^2+1},R^{\mu+\nu}]&=R^{2\mu-(\mu+\nu)\operatorname{ind}(g^2+1)}[r,R^{\mu+\nu}]\\
R^{3\mu}[r^{g^3+1},R^{\mu+\nu}]&=R^{3\mu-(\mu+\nu)\operatorname{ind}(g^3+1)}[r.R^{\mu+\nu}]\\
&\text{etc.}
\end{aligned}
\]

\Needspace{9\baselineskip}
From this we obtain
\begingroup\postdisplaypenalty=10000
\[
\text{I.}\qquad\Omega=[r,R^{\mu+\nu}]\times\Sigma R^{\mu\operatorname{ind}x-(\mu+\nu)\operatorname{ind}(x+1)}
\]
\endgroup
if for $x$ we successively substitute the values $1,g,g^2,g^3\dotsfour g^{n-2}$, except for $g^{\frac12(n-1)}$, or, which evidently amounts to the same thing, if for $x$ we substitute the values $1,2,3,4\dotsfour n-2$, since these values are congruent to those modulo $n$, though their order is changed.

Taking the integer $y$ inverse to $x$ modulo $n$, that is, determined so that $xy\equiv1\pmod n$, we have $\operatorname{ind}x\equiv-\operatorname{ind}y\pmod{n-1}$ and $\operatorname{ind}(x+1)+\operatorname{ind}y\equiv\operatorname{ind}(xy+y)\equiv\operatorname{ind}(1+y)\pmod{n-1}$; hence
\[
\begin{aligned}
\mu\operatorname{ind}x-(\mu+\nu)\operatorname{ind}(x+1)&\equiv-\mu\operatorname{ind}y-(\mu+\nu)\{\operatorname{ind}(y+1)-\operatorname{ind}y\}\\
&\equiv\nu\operatorname{ind}y-(\mu+\nu)\operatorname{ind}(y+1)
\end{aligned}
\]
Therefore, since the inverses of $1,2,3\dots n-2$ coincide with these same numbers, with only their order changed, we also have
\[
\text{II.}\qquad\Omega=[r,R^{\mu+\nu}]\times\Sigma R^{\nu\operatorname{ind}y-(\mu+\nu)\operatorname{ind}(y+1)}
\]
substituting for $y$ successively the numbers $1,2,3\dots n-2$. The same formula follows directly from I, since the numbers $\mu,\nu$ can evidently be interchanged.

Finally, taking the integer $z$ inverse to $x+1$ modulo
\SourceLeaf{263}
$n$, or $xz+z\equiv1\pmod n$, we have $\operatorname{ind}(1-z)\equiv\operatorname{ind}x+\operatorname{ind}z\pmod{n-1}$, $\operatorname{ind}(x+1)\equiv-\operatorname{ind}z\pmod{n-1}$, and hence
\[
\begin{aligned}
\mu\operatorname{ind}x-(\mu+\nu)\operatorname{ind}(x+1)&\equiv\mu(\operatorname{ind}(1-z)-\operatorname{ind}z)+(\mu+\nu)\operatorname{ind}z\\
&\equiv\mu\operatorname{ind}(1-z)+\nu\operatorname{ind}z
\end{aligned}
\]
Thus, as $x$ runs through the values $1,2,3\dots n-2$, the number $z$ must run through the values $2,3,4\dots n-1$, though in another order, and we obtain the third expression
\begingroup\postdisplaypenalty=10000
\[
\text{III.}\qquad\Omega=[r,R^{\mu+\nu}]\times\Sigma R^{\mu\operatorname{ind}(1-z)+\nu\operatorname{ind}z}
\]
\endgroup
substituting for $z$ successively the values $2,3,4\dots n-1$, or, if you prefer,
\[
\begin{aligned}
\text{IV.}\qquad\Omega&=[r,R^{\mu+\nu}]\times\Sigma R^{\mu\operatorname{ind}(n+1-z)+\nu\operatorname{ind}z}\\
&=[r,R^{\mu+\nu}]\times\Sigma R^{\mu\operatorname{ind}z+\nu\operatorname{ind}(n+1-z)}
\end{aligned}
\]
Since $\operatorname{ind}(1-z)=\frac12(n-1)+\operatorname{ind}(z-1)$, our product can also be represented as follows:
\[
\begin{aligned}
\relax[r,R^\mu]\times[r,R^\nu]&=R^{\frac12(n-1)\mu}\{[1,R^{\mu+\nu}]+[r,R^{\mu+\nu}]\times\Sigma R^{\mu\operatorname{ind}(z-1)+\nu\operatorname{ind}z}\}\\
&=R^{\frac12(n-1)\nu}\{[1,R^{\mu+\nu}]+[r,R^{\mu+\nu}]\times\Sigma R^{\mu\operatorname{ind}z+\nu\operatorname{ind}(z-1)}\}
\end{aligned}
\]
where $z$ is always understood to take the values $2,3,4\dotsfour n-1$.

Moreover, in all these formulas the numbers
\[
\mu\operatorname{ind}x-(\mu+\nu)\operatorname{ind}(x+1),\quad\nu\operatorname{ind}y-(\mu+\nu)\operatorname{ind}(y+1),\quad\mu\operatorname{ind}(1-z)+\nu\operatorname{ind}z
\]
etc., can evidently be replaced by their least residues modulo $\beta$.

If $\mu+\nu\equiv0\pmod\beta$, we have
\[
\begin{aligned}
\relax[r,R^\mu][r,R^\nu]&=(n-1)R^{\frac12(n-1)\mu}\\
&\quad+(r+r^2+r^3+\dots+r^{n-1})\\
&\qquad\times(1+R^\mu+R^{2\mu}+R^{3\mu}+\dots+R^{(n-2)\mu}-R^{\frac12(n-1)\mu})
\end{aligned}
\]

\clearpage
\Article{12}
The product $[1,R^\mu]\times[r,R^\nu]$ becomes, by the theorem of Article~4,
\[
\begin{aligned}
&=[r,R^{\mu+\nu}]+R^\mu[r,R^{\mu+\nu}]+R^{2\mu}[r,R^{\mu+\nu}]+\text{etc.}+R^{(n-2)\mu}[r,R^{\mu+\nu}]\\
&=[r,R^{\mu+\nu}]\times(1+R^\mu+R^{2\mu}+R^{3\mu}+\text{etc.}+R^{(n-2)\mu})\\
&=[r,R^{\mu+\nu}]\times\frac{n-1}\beta(1+R^\mu+R^{2\mu}+R^{3\mu}+\text{etc.}+R^{(\beta-1)\mu})
\end{aligned}
\]

\SourceLeaf{264}
From this the product $[1,R^\mu]\times[1,R^\nu]$ expands into
\[
\frac{n-1}\beta[1,R^{\mu+\nu}]\times(1+R^\mu+R^{2\mu}+R^{3\mu}+\text{etc.}+R^{(\beta-1)\mu})
\]
The general product $[r^m,R^\mu]\pdot[r^{m'},R^{\mu'}]$ can now be found without difficulty. For $[r^m,R^\mu]=R^{-\mu\operatorname{ind}m}[r,R^\mu]$ when the value of $m$ is not divisible by $n$, and $=[1,R^\mu]$ when it is divisible; and since a similar transformation holds for the other factor $[n^{m'},R^{\mu'}]$, the multiplication reduces either to the problem of the preceding article or to the cases considered in this article.
